\section{Radial variation and the two-dimensional base case}
\label{sec:radial}

Fix a dimension $n\geq 2$ and suppose that the sharp filling inequality
fails in a compact CAT(0) space.  Let $Y$, $A$, $\mu=\|A\|$, $m$, and
$c>c_n$ be the rescaled space, cycle, mass measure, mass, and constant
provided by Lemma~\ref{lem:extremizer}.  Thus
\eqref{eq:extremizer} and \eqref{eq:normalization} hold, with
$q_0=(n+1)/n$.  The differential notation throughout this section is
that of Proposition~\ref{prop:chart-mass}; in particular, norms and inner
products of spatial covectors use the inverse chart metric.

\begin{proposition}[Radial first variation]
\label{prop:radial}
For every center $y\in Y$, write $r(x)=d(y,x)$.  Every nonnegative
Lipschitz function $g:Y\to\R$ satisfies
\begin{equation}
 \int_Y \bigl(ng+r\,Dr\cdot Dg\bigr)\dd\mu
 \leq n\int_Y rg\sqrt{1-|Dr|^2}\dd\mu.
 \label{eq:radial}
\end{equation}
The square root is defined $\mu$-almost everywhere by
$|Dr|\leq 1$.
\end{proposition}

\begin{proof}
Fix $y$ and $g$.  Put $R=\diam Y$ and $M_g=\|g\|_\infty$, and choose
$h>0$ small enough that $hM_g\leq 1/2$.  For $0\leq t\leq h$, let
\[
 \lambda(t,x)=1-tg(x),\qquad
 F(t,x)=[y,x]_{\lambda(t,x)},\qquad F_t(x)=F(t,x),
\]
where $[y,x]_a$ denotes the point at fraction $a$ along the geodesic
from $y$ to $x$.  In particular, $F_0$ is the identity and
$1/2\leq\lambda\leq1$.  Equal-fraction CAT(0) comparison and the
constant-speed parametrization of a geodesic give
\begin{align*}
 d(F(t,x),F(t',x'))
 &\leq d(x,x')+R\,|tg(x)-t'g(x')|\\
 &\leq \bigl(1+hR\Lip(g)\bigr)d(x,x')
       +RM_g|t-t'|.
\end{align*}
Thus $F$ is jointly Lipschitz.  The currents
\[
 B_h=(F_h)_\#A\in\I_n(Y),\qquad
 D_h=F_\#([0,h]\times A)\in\I_{n+1}(Y)
\]
are integral by Theorem~\ref{thm:current-background}.  Their boundaries
are $\partial B_h=0$ and $\partial D_h=B_h-A$.

\emph{Comparison forms.}
We first establish precise differential upper bounds for these maps.
For $x,x'\in Y$ and two fractions $\lambda,\lambda'\in[0,1]$, Euclidean
comparison for the triangle with vertices $y,x,x'$ gives
\begin{equation}
 d([y,x]_\lambda,[y,x']_{\lambda'})^2
 \leq (\lambda r-\lambda'r')^2
       +\lambda\lambda'\bigl(d(x,x')^2-(r-r')^2\bigr),
 \label{eq:radial-two-point}
\end{equation}
where $r=d(y,x)$ and $r'=d(y,x')$.  Indeed, the expression on the
right is exactly the squared distance between the corresponding
points in the Euclidean comparison triangle.  This remains valid for
degenerate triangles.

Take the disjoint charts $\phi_i:E_i\to Y$ from
Proposition~\ref{prop:chart-mass}, with metric matrices $G_i$,
$J_i=\sqrt{\det G_i}$, and signed multiplicities $\theta_i$.
At almost every density point $z\in E_i$, the metric differential
\eqref{eq:chart-metric} and the derivatives of $r\circ\phi_i$ and
$g\circ\phi_i$ exist.  At such a point write
\[
 \alpha=D(r\circ\phi_i)(z),\qquad
 \gamma=D(g\circ\phi_i)(z),
\]
and abbreviate $r=r(\phi_i(z))$, $g=g(\phi_i(z))$, and
$\lambda=1-tg$.  Thus $\alpha$ and $\gamma$ represent $Dr$ and $Dg$
in this chart.  Equation~\eqref{eq:gradient-lipschitz} yields
$|\alpha|_{G_i^{-1}}\leq1$ and
$|\gamma|_{G_i^{-1}}\leq\Lip(g)$.

Set $F_i(t,z)=F(t,\phi_i(z))$.  The quadratic forms
\begin{align}
 P_{\mathrm{tot}}
 &=\lambda^2(G_i-\alpha\otimes\alpha)
   +(\lambda\alpha-tr\gamma-rg\,dt)
       \otimes(\lambda\alpha-tr\gamma-rg\,dt),
       \label{eq:radial-total-form}\\
 P_{\mathrm{sp}}
 &=\lambda^2(G_i-\alpha\otimes\alpha)
   +(\lambda\alpha-tr\gamma)\otimes(\lambda\alpha-tr\gamma)
       \label{eq:radial-spatial-form}
\end{align}
are nonnegative.  In the first expression, the first summand acts
only on spatial vectors.  For increments $(s,w)$ for which
$(t+s,z+w)\in[0,h]\times E_i$, they satisfy
\begin{equation}
 d(F_i(t+s,z+w),F_i(t,z))^2
 \leq P_{\mathrm{tot}}((s,w),(s,w))
       +o(|s|^2+|w|^2).
 \label{eq:radial-differential-upper}
\end{equation}
To see this, the derivative of $(1-tg)r$ is
$\lambda\alpha-tr\gamma-rg\,dt$, while
\[
 d(\phi_i(z+w),\phi_i(z))^2
 -(r(\phi_i(z+w))-r(\phi_i(z)))^2
 =(G_i-\alpha\otimes\alpha)(w,w)+o(|w|^2).
\]
Substitution into \eqref{eq:radial-two-point} proves
\eqref{eq:radial-differential-upper}; changing the factor
$\lambda\lambda'$ to $\lambda^2$ contributes only a higher-order
error.  Taking $s=0$ gives the corresponding upper bound with
$P_{\mathrm{sp}}$ for the fixed-time map.  These statements are
needed separately for each fixed $y$ and $g$; no common exceptional
set for all centers or all functions is required.

\emph{Determinants of the upper forms.}
We next compute the determinants.  At the point under consideration
choose spatial coordinates orthonormal for $G_i$.  In these
coordinates \eqref{eq:radial-spatial-form} becomes
\begin{align*}
 P_{\mathrm{sp}}
 &=\lambda^2 I-\lambda tr(\alpha\gamma^\top+
                         \gamma\alpha^\top)
                    +t^2r^2\gamma\gamma^\top\\
 &=I-t\bigl(2gI+r(\alpha\gamma^\top+
                         \gamma\alpha^\top)\bigr)+O(t^2).
\end{align*}
The derivative at the identity of $H\mapsto\sqrt{\det H}$ is
$H'\mapsto\tfrac12\operatorname{tr}H'$.  Returning to the original
chart coordinates therefore gives
\begin{equation}
 \frac{\sqrt{\det P_{\mathrm{sp}}}}{J_i}
 =1-t\bigl(ng+r\,Dr\cdot Dg\bigr)+O(t^2).
 \label{eq:radial-spatial-jacobian}
\end{equation}
After decreasing $h$ if necessary, the remainder is bounded in
absolute value by $Ct^2$, where $C$ depends only on
$n,R,M_g,$ and $\Lip(g)$.  In particular it is independent of the
chart and the base point: all the matrix entries in the orthonormal
coordinates are controlled by these bounds.

For the full form, put $H=\lambda^2(I-\alpha\alpha^\top)$ in the
same coordinates and $u=\lambda\alpha-tr\gamma$.  Its matrix in
time-first coordinates is
\[
 \begin{pmatrix}0&0\\0&H\end{pmatrix}
 +\begin{pmatrix}-rg\\u\end{pmatrix}
  \begin{pmatrix}-rg&u^\top\end{pmatrix}.
\]
Its determinant is $(rg)^2\det H$.  For positive definite $H$ this
follows directly by subtracting suitable multiples of the first
row and column, or by a block determinant calculation; replacing
$H$ by $H+\varepsilon I$ and letting $\varepsilon\downarrow0$
proves the identity for every nonnegative $H$.  Since
$\det(I-\alpha\alpha^\top)=1-|\alpha|^2$, we obtain the exact
identity
\begin{equation}
 \frac{\sqrt{\det P_{\mathrm{tot}}}}{J_i}
 =rg\lambda^n\sqrt{1-|Dr|^2}.
 \label{eq:radial-total-jacobian}
\end{equation}
Here $r,g,\lambda$ are nonnegative.  In particular, the formula
also covers $rg=0$ and $|Dr|=1$.

Figure~\ref{fig:radial-variation} separates the fixed-time current
from the swept current and illustrates the unequal-fraction comparison.
The forms above are upper bounds for the deformed metric; the exact
identity is the determinant calculation for the comparison form.
\begin{figure}[htbp]
\centering
\begin{tikzpicture}[>=Latex,font=\small]
  \begin{scope}[xshift=0cm]
    \node[anchor=west,font=\bfseries\small] at (-.3,3.4) {(a) Slice and sweep};
    \draw[fill=blue!5,draw=blue!45!black] (0,.2) rectangle (2.6,2.5);
    \foreach \yy in {.2,.66,1.12,1.58,2.04,2.5}
      \draw[blue!30] (0,\yy)--(2.6,\yy);
    \draw[very thick,blue!65!black] (0,1.58)--(2.6,1.58);
    \draw[->] (-.35,.1)--(-.35,2.7) node[above] {$t$};
    \node[left] at (-.35,.2) {$0$};
    \node[left] at (-.35,2.5) {$h$};
    \node at (1.3,2.2) {product};
    \node at (1.3,.9) {$[0,h]\times A$};
    \node[anchor=west] at (2.75,1.58) {$\{t\}\times A$};
    \node[anchor=west] at (0,-1.7) {$\{t\}\times A\mapsto(F_t)_\#A$};
    \node[anchor=west] at (0,-.2) {$B_h=(F_h)_\#A$};
    \node[anchor=west] at (0,-.7) {$D_h=F_\#([0,h]\times A)$};
    \node[anchor=west] at (0,-1.2) {$\partial D_h=B_h-A$};
  \end{scope}
  \begin{scope}[xshift=6.0cm]
    \node[anchor=west,font=\bfseries\small] at (-.1,3.4) {(b) Unequal fractions};
    \coordinate (y) at (0,.15);
    \coordinate (x) at (3.7,.15);
    \coordinate (xp) at (2.4,2.55);
    \coordinate (p) at ($(y)!.73!(x)$);
    \coordinate (pp) at ($(y)!.53!(xp)$);
    \draw[black!60] (y)--(x)--(xp)--cycle;
    \draw[very thick,blue!65!black] (p)--(pp);
    \foreach \pt in {y,x,xp,p,pp} \fill (\pt) circle (1.4pt);
    \node[below left] at (y) {$\bar y$};
    \node[below right] at (x) {$\bar x$};
    \node[above] at (xp) {$\bar x'$};
    \node[below] at (p) {$p$};
    \node[left] at (pp) {$p'$};
    \node[below] at ($(y)!.36!(x)$) {$\lambda r$};
    \node[above left] at ($(y)!.27!(xp)$) {$\lambda'r'$};
    \node[anchor=west] at (0,-.65) {$\lambda=1-tg(x),\quad\lambda'=1-tg(x')$};
    \node[anchor=west] at (0,-1.2) {$d(F_t(x),F_t(x'))\le |p-p'|$};
  \end{scope}
\end{tikzpicture}
\caption{Radial variation in the whole-current product and its Euclidean
comparison triangle. Panel (a) distinguishes a fixed-time slice from
the swept current; the stack is symbolic and makes no regularity or
injectivity assertion about its image in $Y$. In panel (b), the
barred vertices correspond to $y,x,x'$; $p$ and $p'$ lie at fractions
$\lambda$ and $\lambda'$ on their radial sides. The fractions may
differ even at the same time because $g$ varies with the point.
The comparison distance gives the upper bound
\eqref{eq:radial-two-point}. The exact determinant in
\eqref{eq:radial-total-jacobian} belongs to the quadratic upper form,
rather than an asserted exact differential of the deformed map.}
\label{fig:radial-variation}
\end{figure}


\emph{From whole-current products to mass.}
The spatial determinant will bound the endpoint cycle mass; the full
determinant will bound the swept filling mass. We now pass from the
comparison forms to these two current bounds. The chart formula
\eqref{eq:chart-action} gives, for a bounded Lipschitz coefficient $b$
and Lipschitz functions
$\pi_1,\ldots,\pi_n$ on $Y$,
\begin{equation}
 \begin{split}
 B_h(b,\pi_1,\ldots,\pi_n)
 =\sum_i\int_{E_i}&\theta_i(z)b(F_i(h,z))\\[-2pt]
       &{}\cdot\det D_z(\pi_1\circ F_i(h,\cdot),\ldots,
                 \pi_n\circ F_i(h,\cdot))(z)\dd z.
 \end{split}
 \label{eq:radial-pushforward-action}
\end{equation}
For the homotopy current, the corresponding formula is
\begin{equation}
 \begin{split}
 D_h(b,\pi_1,\ldots,\pi_{n+1})
 =\sum_i\int_0^h\int_{E_i}
       &\theta_i(z)b(F_i(t,z))\\[-2pt]
       &{}\cdot\det D_{(t,z)}
          (\pi_1\circ F_i,\ldots,\pi_{n+1}\circ F_i)(t,z)
          \dd z\dd t.
 \end{split}
 \label{eq:radial-product-action}
\end{equation}
To obtain this formula, apply the interval product to the whole
current $A$. Write $\psi_a(t,x)=\pi_a(F(t,x))$ and
$\widetilde b(t,x)=b(F(t,x))$.  Inserting these functions into
\eqref{eq:interval-product-action} expresses the left side as
\[
 \int_0^h\sum_{a=1}^{n+1}(-1)^{a+1}
 A\bigl(\widetilde b_t\,\partial_t\psi_a(t,\cdot),
   \psi_1(t,\cdot),\ldots,\widehat{\psi_a(t,\cdot)},\ldots,
   \psi_{n+1}(t,\cdot)\bigr)\dd t.
\]
The time derivatives admit bounded Borel versions and exist for
$dt\,d\mu$-almost every $(t,x)$.  For a fixed test tuple they have
uniform bounds, as do the spatial chart-covector norms of the
$\psi_a(t,\cdot)$.  Inserting \eqref{eq:chart-action}, every
cofactor in the resulting expansion is bounded by a constant
times $J_i$, by Euclidean determinant bounds.  The sum of the
absolute integrals is consequently bounded by
$C h\sum_i\int_{E_i}|\theta_i|J_i\dd z=C h m$.
Fubini's theorem therefore permits both the insertion and the
interchange of the sum and integrals.  Expansion in the time
column gives exactly \eqref{eq:radial-product-action}.
The separate derivatives agree almost everywhere with the joint
derivatives of the Lipschitz scalar compositions on
$[0,h]\times E_i$.  This proves the formula without assigning a
boundary to any individual chart piece.

The two action formulas are determinant representations of the
finite-mass currents $B_h$ and $D_h$. Their maps have the respective
quadratic majorants $P_{\mathrm{sp}}(h,z)$ and
$P_{\mathrm{tot}}(t,z)$ from
\eqref{eq:radial-differential-upper}. Consider the measures obtained
by pushing forward
\[
 |\theta_i(z)|\sqrt{\det P_{\mathrm{sp}}(h,z)}\dd z
 \quad\hbox{and}\quad
 |\theta_i(z)|\sqrt{\det P_{\mathrm{tot}}(t,z)}\dd t\dd z
\]
under $F_i(h,\cdot)$ and $F_i$, respectively, and summing over $i$.
They are finite: \eqref{eq:radial-spatial-jacobian} bounds the
spatial density by a uniform multiple of $|\theta_i|J_i$, while
\eqref{eq:radial-total-jacobian} bounds the product density by
$RM_g|\theta_i|J_i$. Their integrals are therefore controlled by
$m$ and $h m$, using \eqref{eq:chart-mass}.
Lemma~\ref{lem:quadratic-mass} now bounds the two current mass
measures by these measures. That lemma includes singular quadratic
forms and permits differentiation null sets to depend on the fixed
test tuple.

Combining these bounds with \eqref{eq:chart-mass},
\eqref{eq:radial-spatial-jacobian}, and
\eqref{eq:radial-total-jacobian}, and using $\lambda^n\leq1$,
gives
\begin{align}
 \M(B_h)
 &\leq m-h\int_Y(ng+r\,Dr\cdot Dg)\dd\mu+O(h^2),
       \label{eq:radial-spatial-mass}\\
 \M(D_h)
 &\leq h\int_Y rg\sqrt{1-|Dr|^2}\dd\mu.
       \label{eq:radial-homotopy-mass}
\end{align}
The constant in the first remainder is finite because the
pointwise remainder was uniform and $\mu(Y)=m<\infty$.

\emph{Using maximality.}
The spatial mass decrease and the swept filling mass are now controlled
by the same radial test. We apply the extremizer inequality to compare
them. Set
\[
 I_g=\int_Y(ng+r\,Dr\cdot Dg)\dd\mu,
 \qquad H_g=\int_Y rg\sqrt{1-|Dr|^2}\dd\mu.
\]
Choose a constant $C$ so that
$\M(B_h)\leq U_h:=m-hI_g+Ch^2$ for all sufficiently small
$h>0$.  Since $m>0$, also $U_h>0$ for such $h$.  The current
$-D_h$ fills $A-B_h$, so \eqref{eq:extremizer} and the
monotonicity of $s\mapsto s^{q_0}$ give
\[
 c(m^{q_0}-U_h^{q_0})
 \leq c(m^{q_0}-\M(B_h)^{q_0})
 \leq V(A-B_h)\leq\M(D_h)\leq hH_g.
\]
Dividing by $h$ and letting $h\downarrow0$ yields
$c q_0m^{1/n}I_g\leq H_g$.  The normalization
$c q_0m^{1/n}=1/n$ in \eqref{eq:normalization} is precisely
\eqref{eq:radial}.
\end{proof}

The radial inequality already settles the first dimension of the
induction. An inverse-square cutoff isolates the Euclidean density at
one center and recovers exactly the sphere area $4\pi$.

\begin{proposition}[Sharp filling in cycle dimension two]
\label{prop:base}
Let $Y$ be a compact CAT(0) space.  Every integral two-cycle
$T\in\I_2(Y)$ admits $S\in\I_3(Y)$ with
\[
 \partial S=T,\qquad
 \M(S)\leq c_2\M(T)^{3/2}
          =\frac{\M(T)^{3/2}}{6\sqrt{\pi}}.
\]
\end{proposition}

\begin{proof}
If this inequality failed, Lemma~\ref{lem:extremizer} in dimension
$n=2$ would produce a normalized nonzero cycle $A$ as above with
\begin{equation}
 m<s_2=4\pi.
 \label{eq:base-small-mass}
\end{equation}
Choose a center $y$ for which Lemma~\ref{lem:euclidean-density}
gives
\begin{equation}
 \liminf_{\varepsilon\downarrow0}
 \varepsilon^{-2}\mu(B(y,\varepsilon))\geq\omega_2=\pi.
 \label{eq:base-density}
\end{equation}
Such a center exists because the density statement holds
$\mu$-almost everywhere and $m>0$.

Let $r=d(y,\cdot)$.  For $0<\varepsilon<1$ such that
$\mu\{r=\varepsilon\}=0$, the function
\[
 g_\varepsilon(x)=\max\{r(x),\varepsilon\}^{-2}
\]
is a nonnegative Lipschitz function on $Y$, so it is admissible
in Proposition~\ref{prop:radial}.  On $\{r<\varepsilon\}$ its
gradient is zero and
\[
 2g_\varepsilon+r\,Dr\cdot Dg_\varepsilon
 =2\varepsilon^{-2},\qquad
 2rg_\varepsilon\sqrt{1-|Dr|^2}
 \leq 2\varepsilon^{-1}.
\]
On $\{r>\varepsilon\}$ the chain rule gives
$Dg_\varepsilon=-2r^{-3}Dr$.  Writing
$b=r^{-1}\sqrt{1-|Dr|^2}$ there, the same two integrands are
$2b^2$ and $2b$, respectively.  The separating level has zero
$\mu$-measure by our choice of $\varepsilon$.  Therefore
\eqref{eq:radial} implies
\begin{align}
 (2-2\varepsilon)\varepsilon^{-2}\mu(B(y,\varepsilon))
 &\leq\int_{\{r>\varepsilon\}}(2b-2b^2)\dd\mu\\
 &\leq \frac{m}{2},
 \label{eq:base-cutoff-estimate}
\end{align}
where the last inequality is the elementary bound
$2b-2b^2=\tfrac12-2(b-\tfrac12)^2\leq\tfrac12$.

Only countably many levels of a real function can carry positive
mass under a finite measure, so eligible $\varepsilon$ tend to
zero.  Taking the lower limit of
\eqref{eq:base-cutoff-estimate} along such radii and using
\eqref{eq:base-density} yields $2\pi\leq m/2$, or $m\geq4\pi$.
This contradicts \eqref{eq:base-small-mass}.  The least filling
mass is therefore at most $c_2\M(T)^{3/2}$, and attainment from
Theorem~\ref{thm:current-background} supplies the required
integral filling.  The zero cycle is filled by the zero current.
\end{proof}

The induction is now initialized, with an attained integral filling in
dimension two. In the remaining sections $n>2$, and the only sharp
lower-dimensional input is the theorem in dimension $n-1$, applied to
boundaries of small restrictions of the extremizing $n$-cycle.
